% ============================================================
%  Chapter 2 -- Finite, Countable, and Uncountable Sets, 2.1-2.14
% ============================================================
\rchapter{2}{Basic Topology}

\begin{roadmap}[Why this chapter exists]
Chapter 1 secured a number system with no gaps. That is a statement about
$\R$ alone, and almost nothing in analysis is about $\R$ alone --- limits,
continuity and integration all concern how points sit \emph{near} one another.
This chapter builds the language for nearness, and it builds it once, for a
general metric space, so that $\R$, $\C$, $\R^k$ and later spaces of functions
are all covered by the same theorems.

\medskip
The chapter's centre of gravity is \textbf{compactness} (2.31--2.42). It is the
property that makes ``infinitely many'' behave like ``finitely many,'' and it
is what will later force a continuous function on $[a,b]$ to be bounded (4.15),
to attain its bounds (4.16), and to be uniformly continuous (4.19). If you take
one thing from Chapter 2, take Theorem 2.41.
\end{roadmap}

\begin{roadmap}[How it is organized, and why in that order]
Five sections, and the first is not topology at all.

\medskip
\centering
\begin{tikzpicture}[scale=0.9, transform shape, node distance=4mm]
\node[rmbox, text width=19mm] (a) {\textbf{2.1--2.14}\\counting\\infinite sets};
\node[rmbox, text width=19mm, right=of a] (b) {\textbf{2.15--2.30}\\metric spaces:\\open, closed};
\node[rmbox, text width=19mm, right=of b] (c) {\textbf{2.31--2.42}\\compactness};
\node[rmbox, text width=19mm, right=of c] (d) {\textbf{2.43--2.47}\\perfect,\\connected};
\draw[rmarrow] (a) -- (b); \draw[rmarrow] (b) -- (c); \draw[rmarrow] (c) -- (d);
\end{tikzpicture}

\medskip
\raggedright
\textbf{The first section is not topology.} Countability is set theory, and
Rudin puts it here because he needs it in two specific places: to know that
$\Q$ is countable (used constantly), and to know that $\R$ is not (2.43, and
the Cantor set at 2.44). You could read \S\S2.15--2.47 first and refer back.

\smallskip
\textbf{The definitions are chosen for what they enable, not for what they
describe.} ``Open,'' ``closed,'' ``limit point'' will look like arbitrary
vocabulary on first reading. They are not: each is exactly the hypothesis some
later theorem needs. Watch where each is used rather than trying to intuit it
in advance.

\smallskip
\textbf{Compactness is defined by open covers} (2.32), which is the strangest
definition in the book and the one students resist hardest. Rudin's reason
appears only at 2.41, where he proves that for subsets of $\R^k$ it coincides
with two more familiar conditions. Until then, treat it as a formal property
and watch what it does.

\smallskip
\textbf{Everything is done for a general metric space} where it can be, and for
$\R^k$ where it cannot. That division is deliberate --- 2.41 and 2.42 are
\emph{false} in general metric spaces, and knowing which theorems need $\R^k$
is worth more than knowing the theorems.
\end{roadmap}

\rsection{Finite, Countable, and Uncountable Sets}

\begin{signpost}[What this section is for]
Two facts are needed later and everything here exists to produce them: $\Q$ is
countable (Corollary to 2.13) and the set of $0$--$1$ sequences is not (2.14).
The first makes $\Q$ small enough to enumerate --- spent in Exercises 2.18 and
2.29, in Chapter 3, and wherever a countable dense set is wanted. The second is what makes $\R$
genuinely bigger than $\Q$, and it reappears as 2.43.

\smallskip
The section also fixes notation for functions (2.1--2.2) that the whole book
uses. Read 2.2 carefully even if the rest is familiar: Rudin's $f^{-1}$ is
defined for \emph{sets} and does not presuppose that $f$ has an inverse.
\end{signpost}

\noindent
We begin this section with a definition of the function concept.

\ritem{2.1}{Definition}
Consider two sets $A$ and $B$, whose elements may be any objects whatsoever,
and suppose that with each element $x$ of $A$ there is associated, in some
manner, an element of $B$, which we denote by $f(x)$. Then $f$ is said to be a
\emph{function} from $A$ to $B$ (or a \emph{mapping} of $A$ into $B$). The set
$A$ is called the \emph{domain} of $f$ (we also say $f$ is defined on $A$) and
the elements $f(x)$ are called the \emph{values} of $f$. The set of all values
of $f$ is called the \emph{range} of $f$.\kd{``Associated, in some manner'' is
not a definition, and Rudin knows it --- he is declining to build set theory,
exactly as he declined to build $\Z$ in Chapter 1. One repair is to define a
function as a set of ordered pairs in which each $x$ appears exactly once;
another, Tao's (his 3.3.1), keeps functions primitive and defines them by any
property $P(x,y)$ passing the ``vertical line test,'' with the ordered-pair
encoding relegated to his Exercise 3.5.10. Nothing in this book depends on the
choice.}

\ritem{2.2}{Definition}
Let $A$ and $B$ be two sets and let $f$ be a mapping of $A$ into $B$. If
$E \subset A$, $f(E)$ is defined to be the set of all elements $f(x)$, for
$x \in E$. We call $f(E)$ the \emph{image} of $E$ under $f$. In this notation,
$f(A)$ is the range of $f$. It is clear that $f(A) \subset B$. If $f(A) = B$, we
say that $f$ maps $A$ \emph{onto} $B$. (Note that, according to this usage,
\emph{onto} is more specific than \emph{into}.)

If $E \subset B$, $f^{-1}(E)$ denotes the set of all $x \in A$ such that
$f(x) \in E$. We call $f^{-1}(E)$ the \emph{inverse image} of $E$ under
$f$.\kw{$f^{-1}$ here is an operation on \emph{subsets}, and it is always
defined --- for every $f$, invertible or not. $f^{-1}(E)$ is simply
``everything that lands in $E$,'' possibly empty. Do not read it as an inverse
function. The clash of notation is universal and permanent, and the only
defence is to check whether the argument is a set or a point.} If $y \in B$,
$f^{-1}(y)$ is the set of all $x \in A$ such that $f(x) = y$. If, for each
$y \in B$, $f^{-1}(y)$ consists of at most one element of $A$, then $f$ is said
to be a \emph{1-1} (one-to-one) mapping of $A$ into $B$. This may also be
expressed as follows: $f$ is a 1-1 mapping of $A$ into $B$ provided that
$f(x_1) \ne f(x_2)$ whenever $x_1 \ne x_2$, $x_1 \in A$, $x_2 \in A$.

(The notation $x_1 \ne x_2$ means that $x_1$ and $x_2$ are distinct elements;
otherwise we write $x_1 = x_2$.)

\begin{deeper}[Why inverse images, and not images, are the useful notion]
Nothing yet suggests it, but $f^{-1}$ is the better-behaved of the two
operations, and Chapter 4 will turn on the fact. Inverse images respect every
set operation:
\begin{gather*}
  f^{-1}(E \cup F) = f^{-1}(E)\cup f^{-1}(F),\\
  f^{-1}(E \cap F) = f^{-1}(E)\cap f^{-1}(F),\\
  f^{-1}(E^c) = \bigl(f^{-1}(E)\bigr)^c .
\end{gather*}
Images respect only the first. $f(E \cap F) \subset f(E) \cap f(F)$ can be
strict --- take $f$ constant and $E$, $F$ disjoint --- and images do not
interact with complements at all.

That asymmetry is why Theorem 4.8 characterises continuity by ``$f^{-1}$ of
every open set is open'' rather than by anything about images --- no
image-based version could work, since a constant map is continuous yet its
images are single points, never open. Even the closure relation survives only
one-sidedly for images: Exercise 4.2 proves
$f(\overline{E}) \subset \overline{f(E)}$ and asks for an example where the
inclusion is proper.
\end{deeper}

\ritem{2.3}{Definition}
If there exists a 1-1 mapping of $A$ onto $B$, we say that $A$ and $B$ can be
put in \emph{1-1 correspondence}, or that $A$ and $B$ have the same
\emph{cardinal number}, or, briefly, that $A$ and $B$ are \emph{equivalent},
and we write $A \sim B$. This relation clearly has the following properties:

It is reflexive: $A \sim A$.

It is symmetric: If $A \sim B$, then $B \sim A$.

It is transitive: If $A \sim B$ and $B \sim C$, then $A \sim C$.

Any relation with these three properties is called an \emph{equivalence
relation}.

\ritem{2.4}{Definition}
For any positive integer $n$, let $\N_n$ be the set whose elements are the
integers $1,2,\dots,n$; let $\N$ be the set consisting of all positive
integers. For any set $A$, we say:
\begin{enumerate}[label=(\alph*), leftmargin=2.2em, itemsep=1pt, topsep=3pt]
  \item $A$ is \emph{finite} if $A \sim \N_n$ for some $n$ (the empty set is
        also considered to be finite).
  \item $A$ is \emph{infinite} if $A$ is not finite.
  \item $A$ is \emph{countable} if $A \sim \N$.
  \item $A$ is \emph{uncountable} if $A$ is neither finite nor countable.
  \item $A$ is \emph{at most countable} if $A$ is finite or countable.
\end{enumerate}
Countable sets are sometimes called \emph{enumerable}, or
\emph{denumerable}.\kw{Rudin's ``countable'' means countably \emph{infinite} --
it excludes finite sets. Many other books use ``countable'' for what he calls
``at most countable,'' and then say ``countably infinite'' for his sense. When
you read a theorem elsewhere, check which convention is in force; several of
Rudin's statements become false under the other one.}

For two finite sets $A$ and $B$, we evidently have $A \sim B$ if and only if
$A$ and $B$ contain the same number of elements. For infinite sets, however,
the idea of ``having the same number of elements'' becomes quite vague, whereas
the notion of 1-1 correspondence retains its clarity.

\ritem{2.5}{Example}
Let $\Z$ be the set of all integers. Then $\Z$ is countable. For, consider the
following arrangement of the sets $\Z$ and $\N$:
\[
\begin{array}{lllllllll}
  \Z: & 0, & 1, & -1, & 2, & -2, & 3, & -3, & \dots\\
  \N: & 1, & 2, & 3,  & 4, & 5,  & 6, & 7,  & \dots
\end{array}
\]
We can, in this example, even give an explicit formula for a function $f$ from
$\N$ to $\Z$ which sets up a 1-1 correspondence:
\[
  f(n) = \begin{cases}
    \dfrac{n}{2} & (n \text{ even}),\\[6pt]
    -\dfrac{n-1}{2} & (n \text{ odd}).
  \end{cases}
\]

\begin{center}
\begin{tikzpicture}[x=1mm, y=1mm]
  \foreach \i/\z in {0/$0$,1/$1$,2/$-1$,3/$2$,4/$-2$,5/$3$,6/$-3$}{
    \node[mlbl] (z\i) at (\i*12,9) {\z};
    \node[mlbl] (n\i) at (\i*12,0) {$\the\numexpr\i+1\relax$};
    \draw[thin, -{Stealth[length=3pt]}] (n\i) -- (z\i);
  }
  \node[mlbl] at (86,9) {$\cdots$};
  \node[mlbl] at (86,0) {$\cdots$};
  \node[lbl, anchor=east] at (-4,9) {$\Z$};
  \node[lbl, anchor=east] at (-4,0) {$\N$};
\end{tikzpicture}
\end{center}
\figcap{Counting the integers by alternating sides of $0$. Every integer is
reached exactly once, which is all ``countable'' asks.}

\ritem{2.6}{Remark}
A finite set cannot be equivalent to one of its proper subsets. That this is,
however, possible for infinite sets, is shown by Example 2.5, in which $\N$ is
a proper subset of $\Z$.

In fact, we could replace Definition 2.4(b) by the statement: $A$ is infinite
if $A$ is equivalent to one of its proper subsets.\kd{This is Dedekind's
definition of infinity, and taking it as primary is attractive because it needs
no reference to $\N$ at all. The two definitions agree, but proving they agree
requires a form of the axiom of choice --- which is why Rudin offers the
Dedekind version as an aside rather than adopting it.}

\ritem{2.7}{Definition}
By a \emph{sequence}, we mean a function $f$ defined on the set $\N$ of all
positive integers. If $f(n) = x_n$, for $n \in \N$, it is customary to denote
the sequence $f$ by the symbol $(x_n)_{n=1}^\infty$, or more briefly by
$(x_n)$, or sometimes by $x_1, x_2, x_3, \dots$. The values of $f$, that is,
the elements $x_n$, are called the \emph{terms} of the sequence. If $A$ is a
set and if $x_n \in A$ for all $n \in \N$, then $(x_n)$ is said to be a
\emph{sequence in $A$}, or a \emph{sequence of elements of $A$}.

Note that the terms $x_1, x_2, x_3, \dots$ of a sequence need not be
distinct.\kw{A sequence is not its range. $(1,1,1,\dots)$ is an infinite
sequence whose range is the one-element set $\{1\}$. Rudin flags this because
the next paragraph slides between the two, and because Chapter 3 will
repeatedly distinguish ``the sequence converges'' from statements about the
set of its values.}

Since every countable set is the range of a 1-1 function defined on $\N$, we
may regard every countable set as the range of a sequence of distinct terms.
Speaking more loosely, we may say that the elements of any countable set can be
``arranged in a sequence.''

Sometimes it is convenient to replace $\N$ in this definition by the set of all
nonnegative integers, i.e., to start with $0$ rather than $1$.

\ritem{2.8}{Theorem}
\emph{Every infinite subset of a countable set $A$ is countable.}

\begin{proof}
Suppose $E \subset A$, and $E$ is infinite. Arrange the elements $x$ of $A$ in
a sequence $(x_n)$ of distinct elements. Construct a sequence $(n_k)$ as
follows:

Let $n_1$ be the smallest positive integer such that $x_{n_1} \in E$. Having
chosen $n_1,\dots,n_{k-1}$ $(k = 2,3,4,\dots)$, let $n_k$ be the smallest
integer greater than $n_{k-1}$ such that $x_{n_k} \in E$.\kn{``Smallest'' is
well-ordering again, the principle used silently in 1.1 and in the proof of
1.20(b). Choosing the least admissible index at each stage is what makes the
construction canonical --- no arbitrary choices are made, so no appeal to the
axiom of choice is needed here. Contrast the proof of 2.12.}

Putting $f(k) = x_{n_k}$ $(k = 1,2,3,\dots)$, we obtain a 1-1 correspondence
between $E$ and $\N$.
\end{proof}

The theorem shows that, roughly speaking, countable sets represent the
``smallest'' infinite: No uncountable set can be a subset of a countable set.

\begin{unpack}[What 2.8 rules out]
Read as a statement about size, 2.8 says there is nothing strictly between
finite and countable. An infinite subset of a countable set cannot be a
``smaller infinity'' --- it is countable, full stop.

That has an immediate consequence used everywhere: to prove some set $S$ is
countable it is enough to exhibit a 1-1 map from $S$ \emph{into} $\N$, or into
any countable set, together with the knowledge that $S$ is infinite. You never
have to construct a bijection. Rudin uses exactly this shortcut at the end of
2.12, where he obtains $S \sim T$ for some $T \subset \N$ and then invokes 2.8.
\end{unpack}

\ritem{2.9}{Definition}
Let $A$ and $\Omega$ be sets, and suppose that with each element $\alpha$ of
$A$ there is associated a subset of $\Omega$ which we denote by $E_\alpha$.

The set whose elements are the sets $E_\alpha$ will be denoted by
$\{E_\alpha\}$. Instead of speaking of sets of sets, we shall sometimes speak
of a \emph{collection} of sets, or a \emph{family} of sets.

The \emph{union} of the sets $E_\alpha$ is defined to be the set $S$ such that
$x \in S$ if and only if $x \in E_\alpha$ for at least one $\alpha \in A$. We
use the notation
\begin{equation}
  S = \bigcup_{\alpha \in A} E_\alpha. \tag{2.1}
\end{equation}
If $A$ consists of the integers $1,2,\dots,n$, one usually writes
\begin{equation}
  S = \bigcup_{m=1}^{n} E_m \tag{2.2}
\end{equation}
or
\begin{equation}
  S = E_1 \cup E_2 \cup \cdots \cup E_n. \tag{2.3}
\end{equation}
If $A$ is the set of all positive integers, the usual notation is
\begin{equation}
  S = \bigcup_{m=1}^{\infty} E_m. \tag{2.4}
\end{equation}
The symbol $\infty$ in (2.4) merely indicates that the union of a countable
collection of sets is taken, and should not be confused with the symbols
$+\infty$, $-\infty$, introduced in Definition 1.23.\kw{Worth pausing on. The
$\infty$ in $\bigcup_{m=1}^\infty$ is not a limit and nothing is approached ---
it abbreviates ``over all $m \in \N$,'' and the union is defined outright by
the membership condition, not as the limit of finite unions. Rudin flags the
clash because Chapter 3 will use the same symbol for genuine limits.}

The \emph{intersection} of the sets $E_\alpha$ is defined to be the set $P$
such that $x \in P$ if and only if $x \in E_\alpha$ for every $\alpha \in A$.
We use the notation
\begin{equation}
  P = \bigcap_{\alpha \in A} E_\alpha, \tag{2.5}
\end{equation}
or
\begin{equation}
  P = \bigcap_{m=1}^{n} E_m = E_1 \cap E_2 \cap \cdots \cap E_n, \tag{2.6}
\end{equation}
or
\begin{equation}
  P = \bigcap_{m=1}^{\infty} E_m, \tag{2.7}
\end{equation}
as for unions. If $A \cap B$ is not empty, we say that $A$ and $B$
\emph{intersect}; otherwise they are \emph{disjoint}.

\ritem{2.10}{Examples}
\begin{enumerate}[label=(\alph*), leftmargin=2.2em, itemsep=4pt, topsep=4pt]
  \item Suppose $E_1$ consists of $1,2,3$ and $E_2$ consists of $2,3,4$. Then
        $E_1 \cup E_2$ consists of $1,2,3,4$, whereas $E_1 \cap E_2$ consists
        of $2,3$.
  \item Let $A$ be the set of real numbers $x$ such that $0 < x \le 1$. For
        every $x \in A$, let $E_x$ be the set of real numbers $y$ such that
        $0 < y < x$. Then
        \begin{enumerate}[label=(\roman*), leftmargin=2.2em, itemsep=1pt, topsep=3pt]
          \item $E_x \subset E_z$ if and only if $0 < x \le z \le 1$;
          \item $\bigcup_{x \in A} E_x = E_1$;
          \item $\bigcap_{x \in A} E_x$ is empty.
        \end{enumerate}
        (i) and (ii) are clear. To prove (iii), we note that for every $y > 0$,
        $y \notin E_x$ if $x < y$. Hence $y \notin \bigcap_{x \in A}
        E_x$.\kn{Example (b)(iii) is the one to remember: an intersection of
        nonempty nested sets can be empty. Each $E_x$ contains points, and any
        \emph{finitely} many of them have a common point, yet nothing survives
        all of them. Theorem 2.36 will show that compactness is exactly what
        forbids this, and 2.39 is the version for $k$-cells.}
\end{enumerate}

\ritem{2.11}{Remarks}
Many properties of unions and intersections are quite similar to those of sums
and products; in fact, the words \emph{sum} and \emph{product} were sometimes
used in this connection, and the symbols $\Sigma$ and $\Pi$ were written in
place of $\bigcup$ and $\bigcap$.

The commutative and associative laws are trivial:
\begin{align}
  A \cup B &= B \cup A; \qquad A \cap B = B \cap A. \tag{2.8}\\
  (A \cup B) \cup C &= A \cup (B \cup C); \qquad
  (A \cap B) \cap C = A \cap (B \cap C). \tag{2.9}
\end{align}
Thus the omission of parentheses in (2.3) and (2.6) is justified.

The distributive law also holds:
\begin{equation}
  A \cap (B \cup C) = (A \cap B) \cup (A \cap C). \tag{2.10}
\end{equation}
To prove this, let the left and right members of (2.10) be denoted by $E$ and
$F$, respectively.

Suppose $x \in E$. Then $x \in A$ and $x \in B \cup C$, that is, $x \in B$ or
$x \in C$ (possibly both). Hence $x \in A \cap B$ or $x \in A \cap C$, so that
$x \in F$. Thus $E \subset F$.

Next, suppose $x \in F$. Then $x \in A \cap B$ or $x \in A \cap C$. That is,
$x \in A$, and $x \in B \cup C$. Hence $x \in A \cap (B \cup C)$, so that
$F \subset E$.

It follows that $E = F$.\kn{The template for every set identity in the book:
prove two inclusions, each by chasing a single element. Rudin writes it out
once here and never again --- later identities are left as ``easily
verified.''}

We list a few more relations which are easily verified:
\begin{align}
  A &\subset A \cup B, \tag{2.11}\\
  A \cap B &\subset A. \tag{2.12}
\end{align}
If $\emptyset$ denotes the empty set, then
\begin{equation}
  A \cup \emptyset = A, \qquad A \cap \emptyset = \emptyset. \tag{2.13}
\end{equation}
If $A \subset B$, then
\begin{equation}
  A \cup B = B, \qquad A \cap B = A. \tag{2.14}
\end{equation}

\ritem{2.12}{Theorem}
\emph{Let $\{E_n\}$, $n = 1,2,3,\dots$, be a countable collection of countable
sets, and put}
\begin{equation}
  S = \bigcup_{n=1}^{\infty} E_n. \tag{2.15}
\end{equation}
\emph{Then $S$ is countable.}

\begin{proof}
Let every set $E_n$ be arranged in a sequence $(x_{nk})$, $k = 1,2,3,\dots$,
and consider the infinite array\kw{This opening clause is where the axiom of
countable choice enters, and Rudin does not mention it. Each $E_n$ \emph{can}
be arranged in a sequence, but there is no canonical way to do it, and we are
choosing one arrangement for each of infinitely many $n$ simultaneously.
Without some choice principle the theorem is not provable --- it is consistent
with ZF that $\R$ is a countable union of countable sets. Compare 2.8, where
``smallest index'' made every choice canonical and no such principle was
needed.}
\begin{equation}
\begin{array}{ccccc}
  x_{11} & x_{12} & x_{13} & x_{14} & \cdots\\
  x_{21} & x_{22} & x_{23} & x_{24} & \cdots\\
  x_{31} & x_{32} & x_{33} & x_{34} & \cdots\\
  x_{41} & x_{42} & x_{43} & x_{44} & \cdots\\
  \multicolumn{5}{c}{\cdots\cdots\cdots\cdots\cdots\cdots\cdots\cdots}
\end{array}
\tag{2.16}
\end{equation}
in which the elements of $E_n$ form the $n$th row. The array contains all
elements of $S$. As indicated by the arrows, these elements can be arranged in
a sequence
\begin{equation}
  x_{11};\; x_{21}, x_{12};\; x_{31}, x_{22}, x_{13};\;
  x_{41}, x_{32}, x_{23}, x_{14};\; \dots \tag{2.17}
\end{equation}
If any two of the sets $E_n$ have elements in common, these will appear more
than once in (2.17). Hence there is a subset $T$ of the set of all positive
integers such that $S \sim T$, which shows that $S$ is at most countable
(Theorem 2.8). Since $E_1 \subset S$, and $E_1$ is infinite, $S$ is infinite,
and thus countable.
\end{proof}

\begin{center}
\begin{tikzpicture}[x=1mm, y=1mm]
  \foreach \r in {0,1,2,3}{
    \foreach \c in {0,1,2,3}{
      \node[mlbl] (e\r\c) at (\c*15, -\r*9)
        {$x_{\the\numexpr\r+1\relax\the\numexpr\c+1\relax}$};}
    \node[mlbl] at (62,-\r*9) {$\cdots$};}
  \node[mlbl] at (0,-38) {$\vdots$};
  % the diagonal sweeps
  \foreach \d in {1,2,3}{
    \draw[guide, -{Stealth[length=3.5pt]}]
      (\d*15-3, 3) .. controls (\d*7.5,-\d*4.5) .. (2.5,-\d*9+2.5);}
  \node[lbl, anchor=west] at (70,-9) {sweep the};
  \node[lbl, anchor=west] at (70,-13) {anti-diagonals};
\end{tikzpicture}
\end{center}
\figcap{Array (2.16). Each anti-diagonal is finite, so sweeping them in turn
reaches every entry after finitely many steps --- which is exactly what an
enumeration must do.}

\begin{unpack}[Why the diagonal sweep, and not the obvious orders]
The natural first attempts both fail, and seeing why is the content of the
proof.

\smallskip
\emph{Row by row} never finishes row $1$, so $x_{21}$ is never reached.
\emph{Column by column} fails the same way. What the anti-diagonals have that
rows do not is \textbf{finiteness}: the $d$th anti-diagonal holds exactly $d$
entries, so every entry sits at a finite distance from the start.

\smallskip
That is the whole idea, and it is worth stating in the form used later: a
countable union of \emph{finite} sets is at most countable, and the array
argument works by exhibiting $S$ as such a union --- the anti-diagonals.
\end{unpack}

\ritem{}{Corollary}
\emph{Suppose $A$ is at most countable, and, for every $\alpha \in A$,
$B_\alpha$ is at most countable. Put}
\[ T = \bigcup_{\alpha \in A} B_\alpha. \]
\emph{Then $T$ is at most countable}, for $T$ is equivalent to a subset of
(2.15).

\ritem{2.13}{Theorem}
\emph{Let $A$ be a countable set, and let $B_n$ be the set of all $n$-tuples
$(a_1,\dots,a_n)$, where $a_k \in A$ $(k = 1,\dots,n)$, and the elements
$a_1,\dots,a_n$ need not be distinct. Then $B_n$ is countable.}

\begin{proof}
That $B_1$ is countable is evident, since $B_1 = A$. Suppose $B_{n-1}$ is
countable $(n = 2,3,4,\dots)$. The elements of $B_n$ are of the form
\begin{equation}
  (b,a) \qquad (b \in B_{n-1},\ a \in A). \tag{2.18}
\end{equation}
For every fixed $b$, the set of pairs $(b,a)$ is equivalent to $A$, and hence
countable. Thus $B_n$ is the union of a countable set of countable sets. By
Theorem 2.12, $B_n$ is countable.

The theorem follows by induction.
\end{proof}

\ritem{}{Corollary}
\emph{The set of all rational numbers is countable.}

\begin{proof}
We apply Theorem 2.13, with $n = 2$, noting that every rational $r$ is of the
form $b/a$, where $a$ and $b$ are integers. The set of pairs $(a,b)$, and
therefore the set of fractions $b/a$, is countable.
\end{proof}

\begin{deeper}[The corollary is the reason for the section]
$\Q$ countable is used more often than any other fact in this chapter. It gives
$\R$ a countable dense subset (by 1.20(b)) --- which is what Exercise 2.18 uses
to build a perfect set avoiding every rational, what the structure theorem for
open sets (Exercise 2.29) rests on, and what makes $\R^k$ separable (Exercise
2.22).

Note also the shape of the argument, which recurs: $\Q$ is not obviously
countable, but it is the image of a set of \emph{pairs}, and pairs are handled
by 2.13. Whenever an object is specified by finitely many integers, this
argument shows there are only countably many of them --- which is how Exercise
2.2 gets the algebraic numbers.

The strict inequality $|\Q| < |\R|$ that follows from 2.14 is worth pausing on
for what it says about 1.20(b): between any two reals there is a rational, yet
the rationals are vanishingly rare. Density and abundance are unrelated.
\end{deeper}

In fact, even the set of all algebraic numbers is countable (see Exercise 2.2).
That not all infinite sets are, however, countable, is shown by the next
theorem.

\ritem{2.14}{Theorem}
\emph{Let $A$ be the set of all sequences whose elements are the digits $0$ and
$1$. This set $A$ is uncountable.}

The elements of $A$ are sequences like $1,0,0,1,0,1,1,1,\dots$.

\begin{proof}
Let $E$ be a countable subset of $A$, and let $E$ consist of the sequences
$s_1, s_2, s_3, \dots$. We construct a sequence $s$ as follows. If the $n$th
digit in $s_n$ is $1$, we let the $n$th digit of $s$ be $0$, and vice versa.
Then the sequence $s$ differs from every member of $E$ in at least one place;
hence $s \notin E$. But clearly $s \in A$, so that $E$ is a proper subset of
$A$.

We have shown that every countable subset of $A$ is a proper subset of $A$. It
follows that $A$ is uncountable (for otherwise $A$ would be a proper subset of
$A$, which is absurd).\kn{The last step is where 2.6 is spent. If $A$ were
countable then $A$ itself would be one of the countable subsets $E$, and the
argument just showed every such $E$ is a \emph{proper} subset of $A$ --- so
$A \subsetneq A$. Rudin's phrasing (``every countable subset is proper'')
avoids having to fix an enumeration of $A$ in advance and then derive a
contradiction, which is the more familiar but slightly clumsier route.}
\end{proof}

\begin{center}
\begin{tikzpicture}[x=1mm, y=1mm]
  \def\rowA{1,0,0,1,0}
  \foreach \r/\row in {0/{1,1,0,1,0}, 1/{0,1,1,0,1}, 2/{1,0,0,0,1},
                       3/{1,1,1,0,0}, 4/{0,0,1,1,1}}{
    \node[lbl, anchor=east] at (-3,-\r*7) {$s_{\the\numexpr\r+1\relax}$};
    \foreach \c [count=\i from 0] in \row {
      \node[mlbl] at (\i*10,-\r*7) {$\c$};}
    \node[mlbl] at (54,-\r*7) {$\cdots$};}
  % highlight the diagonal
  \foreach \d in {0,1,2,3,4}{
    \draw[thin] (\d*10-3.4,-\d*7-3.4) rectangle (\d*10+3.4,-\d*7+3.4);}
  % the new sequence
  \node[lbl, anchor=east] at (-3,-40) {$s$};
  \foreach \c [count=\i from 0] in {0,0,1,1,0}{
    \node[mlbl] at (\i*10,-40) {$\c$};}
  \node[mlbl] at (54,-40) {$\cdots$};
  \draw[guide] (-6,-36) -- (58,-36);
  \node[lbl, anchor=west] at (62,-20) {flip the boxed};
  \node[lbl, anchor=west] at (62,-24) {diagonal};
\end{tikzpicture}
\end{center}
\figcap{Cantor's diagonal process. The new sequence $s$ disagrees with $s_n$ in
place $n$, so it is not in the list --- whatever the list was.}

The idea of the above proof was first used by Cantor, and is called
\emph{Cantor's diagonal process}; for, if the sequences $s_1,s_2,s_3,\dots$
are placed in an array like (2.16), it is the elements on the diagonal which
are involved in the construction of the new sequence.

Readers who are familiar with the binary representation of the real numbers
(base $2$ instead of $10$) will notice that Theorem 2.14 implies that the set
of all real numbers is uncountable. We shall give a second proof of this fact
in Theorem 2.43.

\begin{pitfall}[The binary argument needs one repair]
The map from $0$--$1$ sequences to reals in $[0,1]$ is \emph{not} a bijection:
$0.0111\dots$ and $0.1000\dots$ name the same number, exactly as $0.999\dots =
1$ in base ten. So 2.14 does not immediately transfer.

The repair is cheap. Only the dyadic rationals have two expansions, and there
are countably many of them, so the map fails to be injective on a countable
set --- not enough to make an uncountable set countable. But the gap is real,
and it is the same delicacy that made Rudin warn ``we shall never use
decimals'' at 1.22 and made Tao call the decimal construction of $\R$ the most
complicated of the three.

Theorem 2.43 avoids the issue entirely, which is presumably why Rudin promises
a second proof.
\end{pitfall}
